\documentclass{article}
\usepackage{amsmath, amssymb, amsthm}
\usepackage{geometry}
\geometry{a4paper, margin=1in}
\usepackage{fancyhdr}
\usepackage{hyperref}
\usepackage[english, russian]{babel}
\usepackage[utf8]{inputenc}
\usepackage{listings}
\usepackage{xcolor}

\pagestyle{fancy}
\fancyhf{}
\fancyhead[L]{Algorithms}
\fancyhead[R]{\thepage}

\title{Algorithms 1 \\ Homework 7}
\author{Nikita Zagainov}
\date{\selectlanguage{english}\today}

\lstset{
    language=Python,
    basicstyle=\ttfamily\small,
    keywordstyle=\color{blue},
    stringstyle=\color{red},
    commentstyle=\color{gray},
    backgroundcolor=\color{lightgray!20},
    frame=single,
    rulecolor=\color{black},
    numbers=left,
    numberstyle=\tiny\color{gray},
    stepnumber=1,
    numbersep=10pt,
    showspaces=false,
    showstringspaces=false,
    breaklines=true,
    breakatwhitespace=true,
    tabsize=4,
    captionpos=b
}

\begin{document}

\maketitle

\begin{center}
	\section*{Расширенный алгоритм Евклида. Куча. \\
	  Разные задачи. Графы.}
\end{center}
\begin{enumerate}
	\item
	      Запишем общую форму уравнения:
	      \[
		      ax + by = c
	      \]
	      Пусть нам дана пара целых чисел $x_0, y_0$,
	      удовлетворяющих данному уравнению. Тогда решением
	      уравнения также будут
	      \[
		      x = x_0 + k \cdot \frac{b}{\gcd(a, b)},
		      \quad y = y_0 - k \cdot \frac{a}{\gcd(a, b)}, \quad \forall k \in \mathbb{Z}.
	      \]

	      Для нахождения $x_0, y_0$ можно воспользоваться
	      расширенным алгоритмом Евклида:
	      \begin{itemize}
		      \item Найти $x_0^{\prime}, y_0^{\prime}$,
		            удовлетворяющие уравнению
		            $ax_0^{\prime} + by_0^{\prime} =
			            \gcd(a, b)$.
		      \item Если $\frac{c}{\gcd(a, b)}$
		            не является целым числом,
		            то уравнение не имеет решений.
		            В противном случае
		            берем $x_0 = x_0^{\prime}
			            \cdot \frac{c}{\gcd(a, b)}$ и
		            $y_0 = y_0^{\prime} \cdot
			            \frac{c}{\gcd(a, b)}$.
	      \end{itemize}

	      Алгоритм:
	      \begin{lstlisting}
def extended_gcd(a, b):
    if b == 0:
        return (a, 1, 0)
    d, x, y = extended_gcd(b, a % b)
    return (d, y, x - (a // b) * y)


def solve(a, b, c):
    d, x_0, y_0 = extended_gcd(a, b)
    if c % d != 0:
        return None
    return (x_0 * (c // d), y_0 * (c // d), b // d, a // d)
        \end{lstlisting}

	      Применяя его, получаем:
	      \begin{enumerate}
		      \item[1.] $x = -399 + 55k, \quad y = 247 - 34k$
		      \item[2.] Нет решений
	      \end{enumerate}

	\item Задача $ax + b \equiv 0 \pmod m$ эквивалентна
	      задаче $ax + my = -b$.
	      Эта задача решается аналогично предыдущей.

	      Ответ: $x = 40 + 33k, \quad \forall k \in \mathbb{Z}$.

	\item Эта задача может быть переформулирована следующим
	      образом:
	      \[
		      ax \equiv 1 \pmod m
	      \]
	      Она эквивалентна задаче $ax + my = 1$.
	      По аналогии с
	      первой задачей, решаем задачу
	      $7x + 102y = 1$ и получаем ответ:
	      $x = -29 + 102k, \quad \forall k \in \mathbb{Z}$.

	\item
	      Чтобы добавить элемент в кучу, можно сначала
	      добавить его в конец списка, а затем восстановить
	      свойство кучи. Это можно сделать, сравнивая
	      этот элемент с родителем и меняя их местами,
	      если родитель меньше ребенка. Продолжаем этот
	      процесс, пока родитель не станет больше ребенка,
	      или пока не дойдем до корня кучи.

	      Алгоритм:
	      \begin{lstlisting}
def heappush(heap: list[float], item: float) -> None:
    heap.append(item)

    pos = len(heap) - 1
    while pos > 0:
        parent_pos = (pos - 1) // 2
        if heap[parent_pos] < heap[pos]:
            heap[parent_pos], heap[pos] = heap[pos], heap[parent_pos]
            pos = parent_pos
        else:
            break
          \end{lstlisting}

	      Сложность данного алгоритма --- $O(\log n)$, так как
	      в худшем случае мы пройдем по высоте дерева, которая равна
	      $\lceil \log_2 n \rceil$.

	\item Реализовать очередь с помощью двух стеков можно
	      следующим образом:
	      \begin{itemize}
		      \item \textsc{heapify} сортирует массив за $O(n \log n)$
		            и поочередно складывает их в один стек.
		      \item \textsc{pop} берет верхний элемент из стека за $O(1)$.

		      \item \textsc{push} создаёт временный стек, в который
		            кладёт элементы с верхушки основного стека,
		            пока не дойдёт до элемента, который больше
		            добавляемого. Затем добавляет элемент в
		            основной стек и возвращает элементы из
		            временного стека обратно в основной.
		            Сложность --- $O(n)$.

		      \item \textsc{peek} возвращает верхний
		            элемент стека за $O(1)$.
	      \end{itemize}

	      Реализация:

	      \begin{lstlisting}
class MinQueueStack:
    def __init__(self, elems: list[float]):
        self.stack = sorted(elems, reverse=True)

    def push(self, elem: float):
        temp_stack = []
        while len(self) > 0 and self.peek() < elem:
            temp_stack.append(self.pop())
        self.stack.append(elem)
        while len(temp_stack) > 0:
            self.stack.append(temp_stack.pop())

    def pop(self) -> float:
        return self.stack.pop(-1)

    def peek(self) -> float:
        return self.stack[-1]

    def __len__(self) -> int:
        return len(self.stack)
          \end{lstlisting}

	\item Количество уникальных лесов обхода у графа
	      пути равно количеству способов выбрать подмножество
	      из всех $\lvert V \rvert - 1$ рёбер,
	      то есть $2^{\lvert V \rvert - 1}$.

	      Для любого такого подмножества найдётся
	      ровно один лес обхода. Для начала возьмём
	      дерево, включающее в себя конечную
	      вершину начнём обход с корня этого дерева.
	      Далее аналогично обходим оставшиеся деревья.

	\item
	      Чтобы найти путь длины $n - 1$ в графе, можно
	      использовать \textsc{DFS} от каждой из
	      вершин, пока не найдём искомый путь.
	      Обход графа занимает $O(V + E)$, таких
	      обходов будет $V$, так что общая сложность
	      алгоритма --- $O(V^2 + VE)$.

	      Алгоритм:
	      \begin{lstlisting}
def dfs(graph, node, visited, path):
    visited[node] = True
    path.append(node)
    if len(path) == len(graph):
        return path
    for neighbor in graph[node]:
        if not visited[neighbor]:
            new_path = dfs(graph, neighbor, visited, path)
            if new_path:
                return new_path
    visited[node] = False
    path.pop()
    return None


def find_path(graph: dict):
    visited = {node: False for node in graph}
    for node in graph:
        path = dfs(graph, node, visited, [])
        if path:
            return path
    return None
        \end{lstlisting}
	      Доказать, что путь всегда существует, не получилось :(

	\item
	      \begin{itemize}
		      \item Шаг: 1, вершины: $[a]$
		      \item Шаг: 2, вершины: $[a, b]$
		      \item Шаг: 3, вершины: $[a, b, c]$
		      \item Шаг: 4, вершины: $[a, b, c, f]$
		      \item Шаг: 5, вершины: $[a, b, c, f, e]$
		      \item Шаг: 6, вершины: $[a, b, c, f, e, d]$
	      \end{itemize}

	\item Данная функция возрастает на множестве
	      натуральных чисел, так как все коэффициенты
	      неотрицательны. Поэтому можно воспользоваться
	      бинарным поиском для определения $x$. Если
	      в процессе был найден подходящий $x$, то
	      возвращаем $True$, иначе $False$, так как решение
	      будет нецелым. Для этого нужно найти границы для
	      бинарного поиска. Это можно сделать путём итеративного
	      экспоненциального увеличения границ, пока
	      максимум не превзойдёт $y$.

	      Алгоритм:
	      \begin{lstlisting}
def p(a: list[int], x: int) -> int:
    result = 0
    power = 1
    for i in range(len(a)):
        result += a[i] * power
        power *= x

    return result


def solve(a: list[int], y: int) -> bool:
    if y == a[0]:
        return True
    if y < a[0]:
        return False

    x_min = 0
    x_max = 1

    while p(a, x_max) < y:
        x_min = x_max
        x_max *= 2

    while x_min <= x_max:
        x = (x_min + x_max) // 2
        value = p(a, x)
        if value == y:
            return True
        if value < y:
            x_min = x + 1
        else:
            x_max = x - 1

    return False 
          \end{lstlisting}

	      Заметим, что $x = O({(\frac{y}{a_n})}^{\frac{1}{n}})$,
	      Тогда сложность алгоритма:
	      \[
		      O \left( n {(\frac{y}{a_n})}^{\frac{1}{n}} +
		      n \log \left( {(\frac{y}{a_n})}^{\frac{1}{n}} \right) \right)
	      \]

	\item Данную задачу можно решить следующим образом:
	      \begin{itemize}
		      \item Разбить все этажи на несколько отрезков
		      \item С последнего этажа каждого отрезка
		            сбрасывать шарик, пока он не разобьется.
		            Если уцелел, то переходим к следующему отрезку.
		      \item Проходим все этажи этого отрезка снизу
		            вверх, пока второй шарик не разобьется.
		      \item ответом будет последний этаж, с которого
		            шарик не разбился.
	      \end{itemize}
	      Если количество этажей на отрезке, $k = const$,
	      то сложность алгоритма будет
	      $O(k + \frac{n}{k})$, где $n$ --- количество этажей.
	      Это выражение минимально при $k = \sqrt{n}$:
	      \[
		      O(\sqrt{n} + \sqrt{n}) = O(\sqrt{n})
	      \]
	      Для более оптимального числа бросков, можно
	      изменять $k$ в зависимости от количества оставшихся
	      этажей, выбирая $k = \lceil \sqrt{n} \rceil$, где $n$ ---
	      количество оставшихся этажей.
	      Тогда в худшем случае, когда прочность будет
	      эквивалентна последнему этажу на каком-либо отрезке,
	      надо будет сделать не более 17 бросков.
	      Выбор $k = \lceil \sqrt(n) \rceil$ разобьёт
	      этажи на следующие отрезки:

	      \[
		      \begin{aligned}
			       & [1..10], [11..20], [21..29], [30..38], [39..46],  \\
			       & [47..54], [55..61], [62..68], [69..74], [75..80], \\
			       & [81..85], [86..89], [90..93], [94..96], [97..98], \\
			       & [99..100]
		      \end{aligned}
	      \]
	      Для каждого из этих отрезков сумма количества
	      предшествующих отрезков и длины не превышает 18, но
	      в во время второго прохода можно не кидать шарик
	      с последнего этажа отрезка, так что количество бросков
	      не превышает 17.

\end{enumerate}

\end{document}


